\sectionkicker{Source-backed technical notes}
\subsection*{規模を支える設計――Switch、ZeRO、GLaM、Gopher: Technical Notes}
本欄は記事本文で圧縮した一次資料上の情報を、比較・再検証しやすい形で整理したものである。時系列、確認済みの主張、留意点、一次資料URLを掲載する。
\smallskip
\noindent{\small\textit{Technical Notesでは、一次資料から確認した公開・提供・研究上の事実を記載する。数値・能力評価や提供元・著者による比較表現は、独立再現済みの結果ではなく帰属claimとして扱う。}}\par
\medskip
\subsection*{Theme at a glance}
\addcontentsline{toc}{subsection}{Theme at a glance}
\begin{center}
\footnotesize
\begin{tabularx}{\linewidth}{@{}p{0.20\linewidth}p{0.10\linewidth}p{0.14\linewidth}X@{}}
\toprule
資料 & 位置づけ & 種別 & 時系列 \\
\midrule
Switch Transformers & 主要資料 & 論文 & 2021-01-11 \\
ZeRO-Infinity & 主要資料 & 論文 & 2021-04-16 \\
GLaM & 主要資料 & 論文 & 2021-12-13 \\
Gopher & 主要資料 & モデル & 2021-12-08 \\
Megatron-Turing NLG 530B & 補足資料 & モデル & 2021-10-11 \\
\bottomrule
\end{tabularx}
\end{center}
\normalsize

\begin{technicalnote}{Switch Transformers}{主要資料}
\begingroup
% reader-facing Technical Notes break policy
\widowpenalty=10000
\clubpenalty=10000
\displaywidowpenalty=10000
\begin{tabularx}{\linewidth}{@{}>{\bfseries}p{0.22\linewidth}X@{}}
組織 & Google / authors \\
種別 & 論文 \\
時系列 & 2021-01-11 \\
\end{tabularx}
\smallskip
{\bfseries 一次資料から整理したtechnical points}
\begin{itemize}[leftmargin=1.5em,itemsep=0.35em]
\item \textbf{一次情報で確認できる事実}: Google / authorsの選定済み一次資料では、Switch TransformersはMixture-of-Experts layerで各tokenを単一expertへroutingする簡略化されたsparse expert構成を採り、計算量を抑えながらmodel capacityを拡張する設計を検証した。 著者らはT5系のpretraining・fine-tuning設定でdense modelとの比較、expert数やcapacity factorのablation、training stabilityを評価した。
\end{itemize}
\begin{minipage}{\linewidth}
% reader-facing Technical Notes generic boundary/source tail group
\begin{samepage}
{\bfseries 一次資料}
\begingroup\sloppy
\Urlmuskip=0mu plus 2mu\relax
\begin{itemize}[leftmargin=1.5em,itemsep=0.25em]
\item {\scriptsize\url{http://arxiv.org/abs/2101.03961v3}}
\end{itemize}
\endgroup
\end{samepage}
\end{minipage}
\endgroup
\end{technicalnote}
\medskip

\begin{technicalnote}{ZeRO-Infinity}{主要資料}
\begingroup
% reader-facing Technical Notes break policy
\widowpenalty=10000
\clubpenalty=10000
\displaywidowpenalty=10000
\begin{tabularx}{\linewidth}{@{}>{\bfseries}p{0.22\linewidth}X@{}}
組織 & Microsoft / authors \\
種別 & 論文 \\
時系列 & 2021-04-16 \\
\end{tabularx}
\smallskip
{\bfseries 一次資料から整理したtechnical points}
\begin{itemize}[leftmargin=1.5em,itemsep=0.35em]
\item \textbf{一次情報で確認できる事実}: Microsoft / authorsの選定済み一次資料では、ZeRO-InfinityはGPUだけでなくCPUとNVMe memoryまで利用するheterogeneous memory systemとして、large model trainingのmemory wallへ対処する設計を提案した。 論文はmodel codeのrefactoringを要求せずにmemoryを階層的に利用することを狙い、著者らは512 NVIDIA V100 GPUでのthroughput/scalabilityやsingle DGX-2 nodeでのlarge-model fine-tuning構成を報告した。
\end{itemize}
\begin{minipage}{\linewidth}
% reader-facing Technical Notes generic boundary/source tail group
\begin{samepage}
{\bfseries 一次資料}
\begingroup\sloppy
\Urlmuskip=0mu plus 2mu\relax
\begin{itemize}[leftmargin=1.5em,itemsep=0.25em]
\item {\scriptsize\url{http://arxiv.org/abs/2104.07857v1}}
\end{itemize}
\endgroup
\end{samepage}
\end{minipage}
\endgroup
\end{technicalnote}
\medskip

\begin{technicalnote}{GLaM}{主要資料}
\begingroup
% reader-facing Technical Notes break policy
\widowpenalty=10000
\clubpenalty=10000
\displaywidowpenalty=10000
\begin{tabularx}{\linewidth}{@{}>{\bfseries}p{0.22\linewidth}X@{}}
組織 & Google / authors \\
種別 & 論文 \\
時系列 & 2021-12-13 \\
\end{tabularx}
\smallskip
{\bfseries 一次資料から整理したtechnical points}
\begin{itemize}[leftmargin=1.5em,itemsep=0.35em]
\item \textbf{一次情報で確認できる事実}: Google / authorsの選定済み一次資料では、GLaMは各tokenに対して一部expertだけをactivateするsparsely activated Mixture-of-Experts language modelとして、総parameter数とtoken当たりのactive computationを分離した。 著者らは1.2T total parameter・約97B active parameterの最大構成を含め、few-shot/one-shot task、training energy、dense GPT-3系baselineとの比較を報告した。
\end{itemize}
\begin{minipage}{\linewidth}
% reader-facing Technical Notes generic boundary/source tail group
\begin{samepage}
{\bfseries 一次資料}
\begingroup\sloppy
\Urlmuskip=0mu plus 2mu\relax
\begin{itemize}[leftmargin=1.5em,itemsep=0.25em]
\item {\scriptsize\url{http://arxiv.org/abs/2112.06905v2}}
\end{itemize}
\endgroup
\end{samepage}
\end{minipage}
\endgroup
\end{technicalnote}
\medskip

\begin{technicalnote}{Gopher}{主要資料}
\begingroup
% reader-facing Technical Notes break policy
\widowpenalty=10000
\clubpenalty=10000
\displaywidowpenalty=10000
\begin{tabularx}{\linewidth}{@{}>{\bfseries}p{0.22\linewidth}X@{}}
組織 & DeepMind \\
種別 & モデル \\
時系列 & 2021-12-08 \\
\end{tabularx}
\smallskip
{\bfseries 一次資料から整理したtechnical points}
\begin{itemize}[leftmargin=1.5em,itemsep=0.35em]
\item \textbf{一次情報で確認できる事実}: DeepMindの選定済み一次資料では、Gopher研究ではTransformer-based language modelを数千万parameterから280B parameterまでscaleし、152 taskにわたって性能変化を比較した。 分析はreading comprehensionやfact-checkingなどscaleで伸びる領域だけでなく、logical/mathematical reasoning、bias、toxicityなど伸び方やriskが異なる領域も対象にした。
\end{itemize}
\begin{minipage}{\linewidth}
% reader-facing Technical Notes generic boundary/source tail group
\begin{samepage}
{\bfseries 一次資料}
\begingroup\sloppy
\Urlmuskip=0mu plus 2mu\relax
\begin{itemize}[leftmargin=1.5em,itemsep=0.25em]
\item {\scriptsize\url{http://arxiv.org/abs/2112.11446v2}}
\item {\scriptsize\url{https://deepmind.google/blog/language-modelling-at-scale-gopher-ethical-considerations-and-retrieval/}}
\end{itemize}
\endgroup
\end{samepage}
\end{minipage}
\endgroup
\end{technicalnote}
\medskip

\begin{technicalnote}{Megatron-Turing NLG 530B}{補足資料}
\begingroup
% reader-facing Technical Notes break policy
\widowpenalty=10000
\clubpenalty=10000
\displaywidowpenalty=10000
\begin{tabularx}{\linewidth}{@{}>{\bfseries}p{0.22\linewidth}X@{}}
組織 & Microsoft / NVIDIA \\
種別 & モデル \\
時系列 & 2021-10-11 \\
\end{tabularx}
\smallskip
{\bfseries 一次資料から整理したtechnical points}
\begin{itemize}[leftmargin=1.5em,itemsep=0.35em]
\item \textbf{一次情報で確認できる事実}: Microsoft / NVIDIAの選定済み一次資料では、MicrosoftとNVIDIAはMegatron-LMとDeepSpeedを組み合わせ、tensor・pipeline・data parallelismを統合した3D parallelismで530B parameterのdense autoregressive language modelをtrainingしたと報告した。 first-party報告ではzero-/one-/few-shot taskを含む評価とlarge-scale training systemの構成を示しており、本誌ではparameter順位ではなくdistributed systems設計の一例として扱う。
\end{itemize}
\begin{minipage}{\linewidth}
% reader-facing Technical Notes generic boundary/source tail group
\begin{samepage}
{\bfseries 一次資料}
\begingroup\sloppy
\Urlmuskip=0mu plus 2mu\relax
\begin{itemize}[leftmargin=1.5em,itemsep=0.25em]
\item {\scriptsize\url{https://www.microsoft.com/en-us/research/blog/using-deepspeed-and-megatron-to-train-megatron-turing-nlg-530b-the-worlds-largest-and-most-powerful-generative-language-model/}}
\end{itemize}
\endgroup
\end{samepage}
\end{minipage}
\endgroup
\end{technicalnote}
\medskip
